\chapter{The Textual Foundations}
\label{ch:textual-foundations}

\epigraph{\itshape The text is a sea; we drink from the surface and
imagine we have understood the depth.}{}

\noindent
\xref{ch:two-failures} argued that Islamic economics has failed to
engage with the metaphysical content of its own tradition.
\xref{ch:levels-of-rigor} introduced the methodological apparatus that
will allow us to engage. Before applying that apparatus to specific
claims in Part III, we must establish what the metaphysical content
\emph{is}. This chapter undertakes the survey.

The chapter is the longest in Part I because the textual base is
genuinely large. The Islamic tradition's discourse on rizq, barakah,
qadar, niyyah, riba, zakat, waqf, and the metaphysics of provision
spans the Quran, the hadith collections, fourteen centuries of
\arb{kalam} debate, the entire Akbarian metaphysical tradition
descending from Ibn Arabi, and an active body of modern scholarship.
The reader who is not familiar with this material needs an orientation;
the reader who is needs a structured catalog. This chapter attempts
both.

The chapter proceeds in five parts.
\xref{sec:quran} surveys the Quranic loci, organized by the eight
canonical claims that Part III will treat.
\xref{sec:hadith} surveys the hadith literature on the same themes.
\xref{sec:kalam} treats the classical \arb{kalam} debates,
particularly on \arb{qadar} and \arb{kasb}, which provide the
classical framework for the modern questions.
\xref{sec:akbarian} surveys the Akbarian metaphysical tradition from
Ibn Arabi to Shah Wali Allah, the conceptual register in which the
metaphysical claims have received their most sustained articulation.
\xref{sec:modern-scholarship} surveys the modern scholarship that
has made this material newly available to the contemporary reader.

The chapter is reference material. The reader is not expected to
absorb every citation on first reading; the chapter is structured so
that the relevant material can be located when subsequent chapters
draw on it. A summary at the end consolidates what the reader needs
to take away.

\section{The Quranic loci}
\label{sec:quran}

The Quran is the foundational textual source of the Islamic tradition
and the locus of its most authoritative teachings. The metaphysical
claims at issue in this book are anchored in specific Quranic
passages, which we survey in this section. The organization follows
the eight canonical claims that Part III will treat, with cross-references
forward to the chapters in which each set of verses receives detailed
analysis.

\subsection{On rizq and its decree}

The Quran is unambiguous that provision is from Allah and is decreed
in advance. The locus classicus is \Quran{11:6}:

\begin{quote}
\textit{And there is no creature on earth but that upon Allah is its
provision (rizq), and He knows its place of dwelling and place of
storage. All is in a clear register.}
\end{quote}

The verse asserts three propositions. First, every living creature
has provision that is upon Allah --- meaning, in the classical
exegesis, that Allah has assumed responsibility for it.\footnote{See
al-Tabari, \textit{Jami` al-Bayan}, tafsir of 11:6; al-Razi,
\textit{Mafatih al-Ghayb}, ad loc.; al-Qurtubi, \textit{al-Jami`
li-Ahkam al-Quran}, ad loc. The classical exegetes converge on the
reading that ``upon Allah'' in this construction (\arb{`ala
Allah}) is a stronger claim than ``from Allah''; it asserts a divine
undertaking, not merely a divine source.} Second, the place and
manner of the provision are known to Allah in advance. Third, all
of this is recorded in a clear register. The verse is the textual
basis for the doctrine that provision is decreed and known.

A second key locus is \Quran{17:31}:

\begin{quote}
\textit{Do not kill your children for fear of poverty. We provide
for them and for you. Indeed, killing them is ever a great sin.}
\end{quote}

The verse, which we will treat in detail in Chapter 11, prohibits
infanticide motivated by economic anxiety. The prohibition presupposes
that the underlying anxiety is misplaced: the parents who fear they
cannot afford another child are mistaken about the economic dynamics.
The Quranic claim, in modern terms, is that the addition of a child
does not in fact reduce per-capita household resources, contrary to
the linear-scaling intuition that motivates the parental fear.

The variant phrasing in \Quran{6:151} reverses the order:

\begin{quote}
\textit{Do not kill your children out of poverty; We will provide
for you and for them.}
\end{quote}

The reversal is not insignificant. \Quran{17:31} addresses parents
who anticipate poverty (``fear of poverty''); \Quran{6:151} addresses
parents already in poverty (``out of poverty''). The two verses
together cover both the anticipated and the actual cases. The
underlying claim is the same: provision is not constrained by the
parental decision.

A third locus is \Quran{29:60}:

\begin{quote}
\textit{And how many a creature carries not its own provision! Allah
provides for it and for you, and He is the Hearing, the Knowing.}
\end{quote}

The verse generalizes the claim beyond humans: the wild animals,
which are unable to plan or save, nevertheless receive their
provision; how much more so for human beings who are commanded to
trust in this same providence. The economic significance of this
verse is the explicit decoupling of provision from the creature's
own provisioning capacity.

Other significant verses on rizq include \Quran{42:27} (variation in
provision is intentional and serves moral purposes), \Quran{34:39}
(charity does not deplete provision), \Quran{2:212} (provision can
exceed reckoning, ``without account''), and \Quran{30:37} (variation
in provision contains signs).

\subsection{On charity and increase}

The Quran asserts repeatedly that charity, far from depleting wealth,
increases it. The clearest statement is in \Quran{2:276}:

\begin{quote}
\textit{Allah destroys riba and gives increase for charities. And
Allah does not love every sinning disbeliever.}
\end{quote}

We treat this verse at length in Chapter 13 (riba) and Chapter 14
(charity). The verse is a parallel statement: riba is destroyed
\emph{and} charity is increased. The two clauses imply that the
underlying dynamics are not those of conventional accounting. By
conventional accounting, charitable giving reduces the donor's wealth
by the amount given; by the Quranic statement, it is increased. The
contradiction must be resolved by understanding the verse as a claim
about real dynamics that diverge from nominal accounting.

A second locus is \Quran{2:261}:

\begin{quote}
\textit{The example of those who spend their wealth in the way of
Allah is like a seed [of grain] which grows seven spikes; in each
spike is a hundred grains. And Allah multiplies for whom He wills.}
\end{quote}

The verse uses the metaphor of agricultural multiplication to assert
that charitable spending produces a return seven hundredfold (seven
spikes times one hundred grains each). The metaphor is striking
because the seven hundredfold return does not appear in any
conventional financial instrument; it asserts a dynamic that operates
on a different scale than that captured by interest-rate compounding.

A third locus is \Quran{34:39}:

\begin{quote}
\textit{Say, ``Indeed, my Lord extends provision for whom He wills of
His servants and restricts for him. But whatever you spend (in His
cause) — He will compensate it; and He is the best of providers.''}
\end{quote}

The verse asserts a compensatory mechanism: spending in Allah's cause
is replaced by Allah. The mechanism is not specified at the verse
level; the verse is a Level 1 statement of a relationship that the
hadith and the secondary tradition will articulate further.

\subsection{On riba and destruction}

The Quranic prohibition of riba is articulated across several
passages with increasing severity. The progression is itself
significant: it suggests that the prohibition was understood
gradually rather than imposed all at once, with the final passage
constituting the most absolute statement.

The earliest passage on riba is \Quran{30:39}:

\begin{quote}
\textit{And whatever you give for interest to increase within the
wealth of people will not increase with Allah. But what you give in
zakat, desiring the countenance of Allah --- those are the
multipliers.}
\end{quote}

The verse establishes the contrast: interest does not produce real
increase before Allah; zakat does. The economic claim is precise:
the apparent increase from interest is not real increase; the
apparent depletion from zakat is, in real terms, multiplication.

The second passage is \Quran{3:130}:

\begin{quote}
\textit{O you who have believed, do not consume usury, doubled and
multiplied, but fear Allah that you may be successful.}
\end{quote}

The phrase ``doubled and multiplied'' (\arb{ad`afan muda`afatan})
references the pre-Islamic practice of compound interest with
escalation upon default. The verse focuses on this particularly
egregious form, but the prohibition is not limited to it.

The third passage is the most extensive, \Quran{2:275-279}:

\begin{quote}
\textit{Those who consume riba cannot stand on the Day of Judgement
except as one stands who is being beaten by Satan into insanity.
That is because they say, ``Trade is just like riba.'' But Allah has
permitted trade and forbidden riba. \dots}

\textit{Allah destroys riba and gives increase for charities. \dots}

\textit{O you who have believed, fear Allah and give up what
remains of riba, if you are believers. And if you do not, then be
informed of a war from Allah and His Messenger. But if you repent,
you may have your principal --- thus you do no wrong, nor are you
wronged.}
\end{quote}

The passage contains the strongest legal-theological condemnation in
the Quran: a ``war from Allah and His Messenger'' is declared against
those who persist in riba. The passage also contains the explicit
distinction from trade (the objection ``trade is just like riba''
is acknowledged and rejected), and the explicit recovery of
principal (``you may have your principal''), which establishes that
the prohibition is on the gratuitous return, not on the underlying
loan.

The Quranic basis for treating riba as destructive of real wealth
(rather than merely as morally prohibited) is \Quran{2:276}, where
the verb \arb{yamhaq} (\textit{destroys}, \textit{wipes out},
\textit{erases}) is used. The classical exegetes are unanimous that
this verb describes a real outcome, not a metaphorical one.

\subsection{On gratitude and multiplication}

The clearest Quranic statement on gratitude is \Quran{14:7}:

\begin{quote}
\textit{And [remember] when your Lord proclaimed, ``If you are
grateful, I will surely increase you (in favor); but if you deny,
indeed, My punishment is severe.''}
\end{quote}

The verse contains a divine declaration --- the verb is
\arb{ta'adhdhana}, ``proclaimed,'' which carries the sense of a
formal undertaking --- of an explicit relationship between gratitude
and increase. The increase is unspecified in scope; the verse uses
\arb{azidannakum} (``I will increase you''), which can refer to
material provision, spiritual reward, or both.

The verse is the textual basis for the claim, examined in Chapter
14, that gratitude practice has measurable economic consequences.
The Quranic statement is that the relationship is causal and
divinely guaranteed; the work of Chapter 14 is to specify the
behavioral mediation chain through which the causation operates.

\subsection{On qadar and decree}

The Quran asserts the universality of divine decree in numerous
passages. \Quran{54:49}:

\begin{quote}
\textit{Indeed, all things We created with predestination
(\arb{qadar}).}
\end{quote}

\Quran{57:22}:

\begin{quote}
\textit{No disaster strikes upon the earth or among yourselves
except that it is in a register before We bring it into being.
Indeed that, for Allah, is easy.}
\end{quote}

\Quran{76:30}:

\begin{quote}
\textit{And you do not will except that Allah wills. Indeed, Allah
is ever Knowing and Wise.}
\end{quote}

\Quran{81:29}:

\begin{quote}
\textit{And you do not will except that Allah wills --- Lord of the
worlds.}
\end{quote}

These verses establish the comprehensive scope of decree: created
things, events, and human will itself are within the scope of
divine determination. The classical theological tension --- how to
reconcile this with the moral responsibility that the same Quran
imposes on humans --- is the subject of the \arb{kalam} debates
treated in \xref{sec:kalam}, and of the modern resolution proposed
in Chapter 16.

\Quran{13:11} is also relevant:

\begin{quote}
\textit{For each one are successive [angels] before and behind him
who protect him by the decree of Allah. Indeed, Allah will not change
the condition of a people until they change what is in themselves.}
\end{quote}

The second clause has been the textual basis for theological
positions that emphasize human agency: while everything occurs by
divine decree, divine decree includes a structural conditionality on
human choice. The verse has been read both as supporting libertarian
positions (the Mu`tazila) and as compatible with deterministic ones
(the Ash`ariyya, with the kasb doctrine). The exegetical disputes
are surveyed in \xref{sec:kalam}.

\subsection{On niyyah and intention}

The Quran does not contain a sustained discussion of niyyah as a
technical concept --- the technical development is in the hadith
literature --- but the Quran does establish the principle that the
inner state of the agent is determinative for the moral status of
acts. \Quran{2:282} is, paradoxically, the longest verse in the
Quran and concerns financial transactions:

\begin{quote}
\textit{O you who have believed, when you contract a debt for a
specified term, write it down. \dots And let neither scribe nor
witness be harmed. And if you do, indeed, it is grave disobedience
in you. \dots}
\end{quote}

The verse establishes the requirement of documentation, witnesses,
and structured terms for debt contracts --- the Quranic
financial-engineering specification, written in the seventh century
and largely consonant with what modern contract law requires. The
verse is treated extensively in Chapter 13 as part of the textual
basis for the prohibition of riba.

A second relevant locus is \Quran{2:264}:

\begin{quote}
\textit{O you who have believed, do not invalidate your charities
with reminders or injury, like one who spends his wealth (only) to
be seen by people and does not believe in Allah and the Last Day.}
\end{quote}

The verse explicitly conditions the value of charity on the
intention with which it is given. Charity given to be seen by people
(rather than for Allah's sake) is invalidated. The Quranic
ontological position is that the same physical act --- giving the
same amount of money to the same recipient --- has different real
content depending on the inner state of the giver. The treatment of
this in Chapter 14 develops it further.

\subsection{On wujud, tawhid, and the structure of being}

The Quranic basis for the Akbarian metaphysical doctrines treated in
Chapter 9 is broad. The first declaration of faith --- \arb{la
ilaha illa Allah}, ``there is no being but Allah'' --- has been read
in the Akbarian tradition as a proposition not just about the
exclusivity of worship but about the exclusivity of being. The
clearest verse is \Quran{2:115}:

\begin{quote}
\textit{And to Allah belongs the east and the west. So wherever you
turn, there is the Face of Allah. Indeed, Allah is all-Encompassing
and Knowing.}
\end{quote}

The verse is the locus classicus for the Akbarian doctrine that all
of created reality is theophanic --- a location of divine
self-disclosure. Ibn Arabi cites this verse repeatedly in the
\textit{Futuhat al-Makkiyyah} as the textual ground for
\arb{wahdat al-wujud}.

\Quran{55:26-27} is also a key locus:

\begin{quote}
\textit{Everyone upon it is in a state of vanishing, and there will
remain the Face of your Lord, the Owner of Majesty and Honor.}
\end{quote}

The verse establishes the contingency of all created things and the
absoluteness of the divine. In Akbarian reading, this is not the
claim that things will (in time) cease to exist; it is the claim
that they have, in their own right, no being at all. Their being is
borrowed from the divine. The doctrine is treated in detail in
Chapter 9.

\subsection{Summary of the Quranic loci}

The Quranic foundation for the metaphysical claims that Part III
will treat is comprehensive. The Quran asserts:

\begin{itemize}[leftmargin=2em]
    \item Provision (rizq) is decreed by Allah for every creature.
    \item Children's provision does not derive from parental
    capacity; the addition of children does not deplete household
    welfare.
    \item Charity does not deplete the giver but increases their
    wealth, in real terms.
    \item Interest (riba) destroys real wealth even as it appears
    to grow nominal wealth.
    \item Gratitude (shukr) is the divinely promised condition for
    increase.
    \item All things, including human will, are within the scope of
    divine decree.
    \item Inner state (intention) is constitutive of the moral and
    economic value of acts.
    \item Created things have no independent being; their existence
    is the self-disclosure of absolute being.
\end{itemize}

These claims are not figurative. They are stated in declarative form,
often with the strongest possible divine attribution (``upon Allah
is its provision''). The work of Part III will be to specify their
formal content and, where possible, derive their consequences.

\section{The hadith tradition}
\label{sec:hadith}

The hadith literature --- the recorded sayings and actions of the
Prophet Muhammad and, in some collections, of his companions ---
extends and refines the Quranic teachings on each of the themes
discussed above. The hadith corpus is vast (the major Sunni
collections alone contain tens of thousands of attested reports),
and a comprehensive survey is beyond the scope of this chapter. We
restrict ourselves to the central reports on each theme.

\subsection{On rizq and provision}

The hadith on rizq is extensive. The most frequently cited report is
from \hadith{Sahih Muslim} (2664):

\begin{quote}
\textit{No one of you will leave this world until his provision
(rizq) reaches him in full and his appointed term (\arb{ajal})
has been completed.}
\end{quote}

The hadith makes two claims simultaneously. First, the lifetime
integral of provision is fixed --- it ``reaches in full'' before
death. Second, the time of death is similarly fixed. Together they
imply that no individual dies in a state of having received less than
their decreed provision; striving cannot increase the integral, only
its temporal distribution. This hadith is the textual basis for the
formal treatment in Chapter 10.

A second key report is from \hadith{Sahih al-Bukhari} (3208), in
which the angels record the provision of each soul before its birth.
The hadith establishes that the decree is not merely chronologically
prior to the soul's life but ontologically prior --- it is part of
the structure within which the soul is created.

A third report, from \hadith{Tirmidhi} (2517), is widely cited:

\begin{quote}
\textit{If you all relied on Allah with the proper trust
(\arb{tawakkul}), He would provide for you as He provides for the
birds: they leave (their nests) hungry in the morning and return
filled in the evening.}
\end{quote}

The hadith does not assert that humans should imitate birds in the
sense of foregoing planning. It asserts that the structural truth
the birds exemplify --- that provision arrives without anxiety to
those who properly entrust themselves --- is also true for humans
and is obscured by the human propensity to anxious calculation. The
report is graded \arb{hasan sahih} by Tirmidhi.

\subsection{On barakah}

The hadith on barakah is particularly rich. A central report is from
\hadith{Sahih al-Bukhari} (2079) and \hadith{Sahih Muslim} (1532):

\begin{quote}
\textit{The two parties to a transaction have the option (of
canceling) so long as they have not separated. If they were truthful
and disclosed all defects, their transaction will be blessed. But if
they lied and concealed defects, the blessing of their transaction
will be wiped out.}
\end{quote}

The hadith establishes a precise causal claim: truthfulness and
disclosure produce barakah; lying and concealment destroy it. The
claim is not that the lying transaction will fail in nominal terms
(it might succeed); the claim is that the barakah --- the multiplier
on real value --- will be removed. The report is the textual basis
for the operational treatment of barakah in Chapter 15.

A second report, from \hadith{Tirmidhi} (1209), grade \arb{hasan
sahih}:

\begin{quote}
\textit{The truthful and trustworthy merchant will be with the
prophets, the truthful, and the martyrs (on the Day of Judgement).}
\end{quote}

The eschatological claim is supported by the worldly consequence: the
truthful merchant's transactions are blessed in this world as well.
The two are not in tension; they are different facets of the same
underlying alignment between the merchant's character and the
metaphysical structure of wealth.

A third report, from \hadith{Sahih al-Bukhari} (660), concerns the
barakah in food:

\begin{quote}
\textit{Eat together and not separately, for the blessing
(\arb{barakah}) is in being together.}
\end{quote}

The hadith asserts that the same physical food, eaten in different
configurations, has different real value. The mechanism is implied
to be the social and dispositional context that surrounds the eating;
the barakah is the multiplier on the physical nutrition.

\subsection{On niyyah}

The opening hadith of \hadith{Sahih al-Bukhari} (1) is the most
famous on this theme:

\begin{quote}
\textit{Actions are but by intentions, and every person will have
only what they intended.}
\end{quote}

The hadith is so central to Islamic legal and ethical thought that
many scholars have organized their entire works around it. The claim
is ontological: the same physical action has different reality
depending on the intention behind it. Two people performing the same
prayer movements with different intentions are doing different
things; two people giving the same charity with different intentions
have given different charities. The Quranic verses on intention
(\Quran{2:264} and others) are systematized by this hadith into a
general principle.

The principle is the textual basis for the formal treatment of
niyyah throughout this book, particularly in Chapters 14 and 15.

\subsection{On gratitude}

A famous report from \hadith{Sahih Muslim} (2999) on gratitude:

\begin{quote}
\textit{The believer's affair is wholly good. If something good
happens, he is grateful, and that is good for him. If something bad
happens, he is patient, and that is good for him. This is for none
but the believer.}
\end{quote}

The hadith establishes the structural significance of gratitude: it
is the disposition that converts a positive event into a multiplied
good for the believer. The mediation chain treated in Chapter 14
articulates the mechanism.

A second report, from \hadith{Tirmidhi} (2374), is widely cited:

\begin{quote}
\textit{Look at those below you in worldly status, not at those
above you. This is more likely to keep you from belittling Allah's
favors upon you.}
\end{quote}

The hadith is a practical instruction for cultivating gratitude. The
psychological mechanism --- comparison-based satisfaction is more
sustainable when the comparison is downward --- is well-established
in modern hedonic psychology. The classical Islamic tradition was
making the recommendation on the basis of its understanding of the
relationship between disposition and well-being; the modern
literature confirms it on independent grounds.

\subsection{On qadar}

The hadith on qadar is extensive and the subject of one of the
canonical chapter divisions in the major collections. A central
report is from \hadith{Sahih Muslim} (8), the so-called Hadith of
Gabriel:

\begin{quote}
\textit{Faith is to believe in Allah, His angels, His books, His
messengers, the Last Day, and to believe in qadar, in its good and
its evil.}
\end{quote}

Belief in qadar --- in both its desirable and undesirable aspects ---
is established here as one of the six articles of faith in Sunni
Islam. The hadith establishes the centrality of the doctrine but
does not adjudicate the question of how to reconcile decree with
human responsibility. That question is left to the \arb{kalam}
tradition, which we survey in \xref{sec:kalam}.

A second important report, from \hadith{Sahih al-Bukhari} (3208):

\begin{quote}
\textit{The creation of one of you is collected in his mother's womb
in forty days as a drop, then he becomes a clot for a similar
period, then he becomes a piece of flesh for a similar period. Then
Allah sends an angel and orders him to write four things: his
provision (rizq), his lifespan (\arb{ajal}), his deeds, and
whether he will be among the wretched or the felicitous. \dots}
\end{quote}

The hadith specifies four things that are decreed and recorded for
the soul before birth: provision, lifespan, deeds, and ultimate
disposition. The decree of deeds, in particular, has been the
hardest theological problem; it is the textual basis for the kalam
disputes treated in the next section.

\subsection{On charity and increase}

A famous report from \hadith{Sahih Muslim} (2588):

\begin{quote}
\textit{Charity does not decrease wealth.}
\end{quote}

The brevity of the hadith belies its importance. The claim is
direct, declarative, and contradicts conventional accounting.
Charity, in conventional accounting, decreases wealth by the amount
given. The hadith asserts that, in reality, it does not. The
treatment in Chapter 14 develops the mechanisms by which this
counterintuitive claim is true.

A second report, from \hadith{Tirmidhi} (664):

\begin{quote}
\textit{Charity extinguishes sins as water extinguishes fire.}
\end{quote}

The hadith makes a different claim --- about the spiritual effect of
charity rather than its material consequences --- but the structural
parallel is significant: charity has a real effect that exceeds and
in fact reverses what conventional accounting would suggest.

\subsection{Summary of the hadith tradition}

The hadith corpus extends and specifies the Quranic teachings in
several important ways. It establishes:

\begin{itemize}[leftmargin=2em]
    \item The pre-natal recording of rizq, ajal, deeds, and
    disposition for every soul.
    \item The lifetime invariance of total provision.
    \item The causal role of truthfulness and disclosure in
    producing barakah; the destructive effect of lying and
    concealment.
    \item The constitutive role of intention for the moral and
    economic content of acts.
    \item The status of qadar as one of the articles of faith.
    \item The non-depletion of wealth by charity.
    \item The structural significance of gratitude as the
    disposition that converts events into goods for the believer.
\end{itemize}

The hadith material is the bridge between the Quranic text and the
later theological development. It provides the specific propositions
that the kalam tradition will systematize and that the Akbarian
tradition will philosophically articulate.

\section{The classical kalam debates}
\label{sec:kalam}

\arb{Kalam} is the discipline of speculative theology in the
Islamic tradition. It emerged in the second Islamic century in
response to internal Muslim debates about the nature of the divine
attributes, the status of the Quran, the status of the grave sinner,
and the question of human freedom and divine decree. By the fourth
century the major schools (Mu`tazila, Ash`ariyya, Maturidiyya, and
others) had crystallized into competing systematic theologies, each
of which gave a comprehensive account of the relationship between
divine attributes, divine action, and human action.

For the present book, the most relevant kalam debate is the debate
on \arb{qadar} and human action. We survey the three main Sunni
positions, then briefly discuss the Mu`tazili position. The dispute
is not merely of historical interest; it provides the classical
framework that the modern reformulation in Chapter 16 must respect.

\subsection{The Mu`tazili position}

The Mu`tazila were the earliest formal kalam school, dominant in the
Abbasid intellectual establishment from roughly the second to the
fifth Islamic centuries. On the question of qadar, the Mu`tazila
adopted a position that emphasized human free will. Their argument
was theological and ethical: divine justice (\arb{`adl}, one of
the five Mu`tazili principles) requires that humans be genuinely
responsible for their actions, and genuine responsibility requires
genuine choice. If God created human actions directly, then God would
be the cause of human sin, which is incompatible with divine
justice. Therefore humans must create their own actions.

The Mu`tazili position thus separates divine creative activity from
human action. God created the human being and the human capacity for
action; the actions themselves are generated by the human. The
divine knowledge of these actions is foreknowledge, but not
causation. God knows what humans will do without making them do it.

The Mu`tazila are sometimes characterized as ``rationalist'' in
contrast to the more ``traditionalist'' Ash`ariyya. The
characterization is partly fair --- the Mu`tazila did emphasize
rational principles such as divine justice as constraints on
admissible theology --- but partly misleading. The Mu`tazila were
working out the implications of Quranic teachings as they understood
them, and their position has its own scriptural support (notably
\Quran{18:29}: ``Whoever wills, let him believe; whoever wills, let
him disbelieve'').

The Mu`tazili school declined institutionally after the fifth
Islamic century, though some of its positions were preserved within
Shi`i Imami theology and continue to be defended in the modern
period.

\subsection{The Ash`ari position: kasb}

The Ash`ariyya, founded by Abu al-Hasan al-Ash`ari in the third
Islamic century, took the opposite starting point. Divine omnipotence
and divine creation of all things require that God be the creator
of human actions as well. There is no event in the universe that is
not directly created by God in the moment of its occurrence. The
Mu`tazili separation of human action from divine creation, on this
view, compromises divine omnipotence.

But Ash`arism faces an obvious problem: if God creates human
actions, what is the moral basis for human responsibility? The
Mu`tazila had identified this problem clearly. The Ash`ari response
is the doctrine of \arb{kasb} (acquisition).

\begin{conceptbox}[The doctrine of kasb]
\arb{Kasb} is the technical term for the human role in actions
that are created by God. Although God creates the action, the human
\emph{acquires} it through some mode of involvement that the
Ash`ari tradition variously describes (volition, contiguity, or
the use of the action). The acquisition is sufficient to make the
human responsible for the action, even though the action itself is
divinely created.
\end{conceptbox}

The doctrine has been criticized, both at the time and since, for
being formally unstable: if the human's role is merely
``acquisition,'' it is unclear what difference this makes
ontologically, and so unclear what the basis of moral responsibility
is. Defenders of the doctrine argue that the criticism misunderstands
the nature of acquisition; what kasb names is a real if difficult-to-articulate
relationship between the human and the divinely-created action.

For our purposes, the Ash`ari position has two significant
implications. First, it commits to the view that the human does not
independently exist as a causally efficacious agent; the human is a
locus of divinely created acts. Second, it requires some additional
ontological category (kasb) to make sense of moral responsibility.
Both of these features will be revisited in Chapter 9, where the
Akbarian tradition's treatment of human agency provides a more
philosophically articulated version of the same insight.

\subsection{The Maturidi synthesis}

The Maturidi school, founded by Abu Mansur al-Maturidi in the third
Islamic century in Samarkand, occupies a middle position between
the Mu`tazili and Ash`ari extremes. The Maturidi doctrine
distinguishes between the \emph{capacity} for action (which is
created by God) and the \emph{exercise} of the capacity (which is
the human's own). When the human exercises the capacity, both the
exercise and the resulting action are real, and the human is
responsible for them; but the entire structure within which the
exercise occurs is divinely created.

The position has been criticized as a less philosophically rigorous
version of the Ash`ari doctrine, but it has the virtue of
distinguishing the levels at which divine and human causality
operate. The capacity is divine; the exercise is human; the action
is the joint product. The doctrine has been the dominant Sunni
theology in the Hanafi tradition, including most of the Indian
subcontinent and Ottoman territory.

\subsection{The unresolved tension and the modern reformulation}

None of the three positions resolves the tension between qadar and
ikhtiyar to universal satisfaction. The Mu`tazili position
preserves human responsibility at the cost of restricting divine
creative activity. The Ash`ari position preserves divine omnipotence
at the cost of making moral responsibility opaque. The Maturidi
position attempts a compromise but inherits criticism from both
sides.

The modern reformulation, which we develop in Chapter 16, does not
resolve the dispute on classical theological terms but dissolves it
by showing that the apparent contradiction between determinism and
the reality of choice was a product of Newtonian assumptions about
what determinism implies. In Bohmian mechanics, particles have fully
determined trajectories \emph{and} real motion; in QBism, probability
is irreducibly perspectival, so divine knowledge can be perfect
while human knowledge is finite without contradiction. Either
formulation provides a structure in which qadar and ikhtiyar can
both be real without conflict.

The modern reformulation is not a refutation of the kalam tradition.
It is a complement: where the classical theologians had to work with
the conceptual apparatus available in the second to the fifth Islamic
centuries, the modern thinker has additional resources. The classical
positions can now be evaluated in light of these resources. The
Ash`ari emphasis on divine creation of all events maps cleanly onto
the Bohmian picture, with the human ``kasb'' corresponding to the
particle's actual trajectory through the configuration space. The
Mu`tazili emphasis on real human agency maps cleanly onto the QBist
picture, with the human's epistemic state being genuinely productive
of action without compromising divine knowledge.

We return to this in Chapter 16.

\section{The Akbarian metaphysical tradition}
\label{sec:akbarian}

The Akbarian tradition is the metaphysical school descending from
Muhyi al-Din ibn al-Arabi (d.~638/1240). The school's central doctrine
is \arb{wahdat al-wujud}, the unity of being, which we treat in
detail in Chapter 9. For the present chapter, we provide an
orientation to the major figures and their contributions.

\subsection{Ibn Arabi (d.~1240)}

Born in Murcia, Andalusia, in 560/1165, Ibn Arabi traveled
extensively across the Muslim world before settling in Damascus,
where he died in 638/1240. His written corpus is enormous; the
\textit{Futuhat al-Makkiyyah} (Meccan Revelations) alone runs to
several thousand pages in standard editions. The \textit{Fusus
al-Hikam} (Bezels of Wisdom), a much shorter and more concentrated
work, is the most studied summary of his metaphysics.

Ibn Arabi's central contributions, for our purposes, are:

\begin{enumerate}[leftmargin=2em]
    \item The doctrine of \arb{wahdat al-wujud}, that there is
    only one being (the divine being) and that all created things
    are loci of its self-disclosure.
    \item The articulation of the divine names (\arb{al-asma'
    al-husna}) as the principles by which absolute being unfolds into
    determinate being.
    \item The doctrine of \arb{tajalli} (theophany,
    self-disclosure) as the active process by which absolute being
    manifests in the loci.
    \item The doctrine of \arb{khayal} (imagination) as a real
    ontological domain, intermediate between pure spirit and pure
    matter.
    \item The doctrine of \arb{insan al-kamil} (the perfect
    human) as the locus that manifests all the divine names in
    equilibrium.
\end{enumerate}

Each of these is developed in Chapter 9. For the present chapter, we
note that the doctrines together form a system in which the
metaphysical claims of the Quranic and hadith material receive their
philosophical articulation. The claim that all things have their
provision from Allah is, in Akbarian vocabulary, the claim that all
things are loci of the divine name \textit{ar-Razzaq} (the
Provider). The claim that wealth carries differential blessing is,
in Akbarian vocabulary, the claim that the loci of the divine names
manifest with differential intensity according to the structure of
the locus. And so on through the eight canonical claims.

\subsection{Mulla Sadra (d.~1641)}

Sadr al-Din Muhammad al-Shirazi, known as Mulla Sadra, was the
greatest Safavid-era Iranian philosopher. His major work, the
\textit{al-Asfar al-arba`a} (The Four Journeys), is a systematic
synthesis of Avicennan philosophy, Suhrawardi's illuminationism, and
the Akbarian metaphysical tradition. Mulla Sadra's original
contribution is the doctrine of \arb{harakat al-jawhariyyah}
(substantial motion), which holds that all created beings are in
constant ontological flux --- not merely accidentally changing, but
substantially becoming.

For our project, Mulla Sadra's relevance is twofold. First, his
doctrine of substantial motion provides a philosophical articulation
of the Akbarian doctrine of \arb{tajalli}: the constant becoming
of created things is the constant self-disclosure of absolute being.
Second, his treatment of the relationship between divine knowledge
and creation --- in which divine knowledge is not added to the
creation but is identical with it --- provides a framework for
understanding qadar that is significantly different from the
classical kalam framework.

Mulla Sadra's tradition continues actively in contemporary Iranian
philosophical scholarship and is increasingly available in English
translation through the work of Sayyid Hossein Nasr and others.

\subsection{Shah Wali Allah Dehlawi (d.~1762)}

Shah Wali Allah, the great eighteenth-century Indian scholar, was
both a hadith scholar of the first rank and a metaphysical thinker
in the Akbarian tradition. His major work, the
\textit{Hujjat Allah al-Baligha} (The Conclusive Argument from God),
is a systematic treatment of the rational basis of Islamic law and
practice that draws extensively on Akbarian ontology to articulate
the metaphysical foundations of legal and economic prescriptions.

For our project, Shah Wali Allah is significant because he is one of
the few figures in the tradition who explicitly connects the
metaphysical doctrines to the social and economic order. His
treatment of zakat, of inheritance, of riba, and of waqf draws on
the Akbarian framework to articulate \emph{why} the Islamic
prescriptions take the form they do. He thus provides a bridge from
the metaphysical tradition to the economic content that this book
is concerned with.

The relevant sections of \textit{Hujjat Allah al-Baligha} are
discussed in Chapter 17, where they are integrated into the
treatment of the four anti-concentration mechanisms.

\subsection{The continuing tradition}

The Akbarian tradition is not a closed historical school. It has
continued in various forms across the Muslim world, particularly in
the Sufi orders that took Ibn Arabi as a major intellectual
influence, in the Ottoman ulama tradition, in the Indian Naqshbandi
and Chishti traditions, and in the contemporary Iranian
philosophical schools that descend from Mulla Sadra. The
contemporary Western academic engagement with the tradition,
discussed in the next section, is itself a product of this living
continuity.

For the present chapter, we note that the Akbarian framework is the
philosophical register in which the metaphysical claims of the
Quranic and hadith material have received their most sustained
articulation. The framework is genuinely available; it is not a
museum piece. The contemporary thinker who wishes to engage with the
metaphysical content of the Islamic tradition has the Akbarian
framework as one of the most rigorous tools available.

\section{The modern scholarship}
\label{sec:modern-scholarship}

A productive modern scholarship has made the Akbarian and broader
Islamic metaphysical traditions newly available to the contemporary
reader. We survey the most significant figures.

\subsection{Sayyid Hossein Nasr (b.~1933)}

Nasr is the leading contemporary expositor of traditional Islamic
philosophy in the Western academic context. His major works ---
\textit{An Introduction to Islamic Cosmological Doctrines} (1964),
\textit{Knowledge and the Sacred} (1981), \textit{Islamic Philosophy
from Its Origin to the Present} (2006), and others --- have
established the framework within which the Akbarian and Sadrian
traditions are studied in the West. Nasr's central claim is that the
metaphysical traditions of Islamic thought are not merely of
historical interest but constitute living philosophical resources
that can address contemporary intellectual problems --- particularly
the loss of any metaphysical orientation in modern Western thought.

For our project, Nasr's significance is twofold. First, his
expository work has made the primary texts and the structure of the
Islamic philosophical tradition accessible to non-specialists.
Second, his theoretical claim --- that traditional Islamic
metaphysics has contemporary intellectual relevance --- is precisely
the claim that this book is undertaking to substantiate at the level
of formal economic theory.

\subsection{William Chittick (b.~1943)}

Chittick, professor at SUNY Stony Brook, is the leading Western
academic interpreter of Ibn Arabi. His major works ---
\textit{The Sufi Path of Knowledge: Ibn al-`Arabi's Metaphysics of
Imagination} (1989), \textit{The Self-Disclosure of God: Principles
of Ibn al-`Arabi's Cosmology} (1998), \textit{Ibn al-`Arabi's
Metaphysics of Imagination} (1989) --- constitute the most
philologically rigorous and philosophically careful exposition of
the Akbarian system available in any Western language.

For our project, Chittick's work is foundational. The treatment in
Chapter 9 draws extensively on his exposition. The lectures
recommended in the preparatory primer accompanying this volume are
all by Chittick.

\subsection{Toshihiko Izutsu (1914--1993)}

The Japanese philosopher and Islamicist Toshihiko Izutsu produced
the most penetrating comparative study of Akbarian metaphysics with
non-Islamic philosophical traditions. His \textit{Sufism and Taoism:
A Comparative Study of Key Philosophical Concepts} (1984) is the
classic comparative treatment, drawing structural correspondences
between Akbarian metaphysics and the Daoist tradition. Izutsu's work
demonstrated that the Akbarian system is comparable to other
sophisticated metaphysical traditions and admits of rigorous
analytical treatment.

For our project, Izutsu's example is methodologically important. He
treated Akbarian thought with the philosophical seriousness that
Western academic philosophy reserves for the Western canon. The
result vindicated the seriousness; the method is the one this book
attempts to extend.

\subsection{Syed Muhammad Naquib al-Attas (b.~1931)}

The Malaysian philosopher al-Attas has developed an original
articulation of Islamic metaphysics that synthesizes the Akbarian
tradition with the Avicennan and Ash`arite traditions and applies
the synthesis to questions of education, psychology, and social
order. His major works include \textit{Prolegomena to the
Metaphysics of Islam} (1995) and \textit{Islam and Secularism}
(1978).

Al-Attas's work is particularly significant for the present project
because of his articulation of the Islamic concept of knowledge
(\arb{`ilm}) and its relation to action, and because of his
extended critique of Western secularism that draws on the Islamic
metaphysical tradition. The framework that al-Attas develops for
distinguishing Islamic from secular epistemology is presupposed in
the methodological discussions of \xref{ch:levels-of-rigor}.

\subsection{Sachiko Murata (b.~1943)}

Murata's \textit{The Tao of Islam: A Sourcebook on Gender Relationships
in Islamic Thought} (1992) is a comprehensive treatment of the
Islamic metaphysical understanding of complementarity and polarity.
Though the book's stated topic is gender, the underlying concern is
the broader Islamic ontological structure of opposites, and the work
is the most accessible introduction to the Akbarian doctrine of the
divine names as articulating the structure of created reality.

For our project, Murata's treatment of the divine names is the
clearest exposition for non-specialists. Chapter 9 draws on it.

\subsection{Other modern scholars}

A comprehensive list would be longer. Other figures whose work bears
on this project include:

\begin{itemize}[leftmargin=2em]
    \item \textbf{Henry Corbin} (1903--1978), whose work on
    Suhrawardi and on the imaginal world (\arb{`alam al-mithal})
    is the foundation of the Western scholarly understanding of the
    Akbarian doctrine of imagination.
    \item \textbf{Annemarie Schimmel} (1922--2003), whose
    \textit{Mystical Dimensions of Islam} (1975) remains the
    standard introduction to Sufi thought broadly.
    \item \textbf{Caner Dagli}, whose contemporary work on Ibn
    Arabi extends Chittick's program.
    \item \textbf{James Morris}, whose translations and commentaries
    on the \textit{Futuhat al-Makkiyyah} have made primary texts
    accessible.
    \item \textbf{Mohammad Hashim Kamali}, whose work on Islamic
    legal theory provides the technical background for the
    integration of metaphysical and legal claims.
    \item \textbf{Wael Hallaq}, whose historical-philosophical work
    on Islamic law illuminates the framework within which the
    economic prescriptions operated.
\end{itemize}

This list is illustrative, not exhaustive. The bibliography of the
book gives a fuller catalog.

\section{Synthesis: what the textual base provides}

The textual base surveyed in this chapter provides four kinds of
material that the rest of the book will use.

First, it provides the \emph{core propositions} that Part III will
treat. The eight canonical claims are not arbitrary selections; each
is anchored in multiple Quranic verses, multiple hadith reports, and
extensive classical commentary. The claims are at the center of the
tradition's economic teaching, not at the periphery.

Second, it provides the \emph{conceptual vocabulary} in which the
claims are articulated. The terms rizq, barakah, qadar, niyyah, riba,
zakat, waqf, wujud, tajalli are not casual descriptors. They are
technical terms with developed meanings within the tradition. The
formal treatments of Part III must respect these meanings; the
treatments cannot redefine the terms to suit modern formal needs.

Third, it provides the \emph{interpretive precedents} that constrain
admissible readings. Fourteen centuries of \arb{tafsir} (Quranic
exegesis), \arb{sharh} (hadith commentary), and theological
debate have established what readings of the core texts are
defensible and what readings are not. The modern formal treatment
must be consistent with the interpretive tradition; it cannot
import readings that the tradition would not recognize.

Fourth, it provides the \emph{philosophical articulation} ---
particularly through the Akbarian tradition --- in which the
metaphysical claims have received their most rigorous classical
treatment. The Akbarian system is the bridge between the textual
material and the formal apparatus that Part II will develop. It
permits the reader to see the metaphysical claims not as isolated
exhortations but as elements of a unified ontology.

The reader who has worked through this chapter has the textual base
needed to follow the rest of the book. Specific verses, hadiths, and
classical positions will be cited as they are needed in the
substantive chapters; the present chapter functions as the reference
to which those citations refer.

\begin{summarybox}[Chapter summary]
\textbf{Quranic loci}: each of the eight canonical metaphysical
claims is anchored in multiple Quranic verses. The verses are
declarative, not figurative, and assert specific propositions about
the dynamics of provision, blessing, decree, intention, interest,
charity, and the structure of being.

\textbf{Hadith tradition}: extends and specifies the Quranic
teachings, particularly on the pre-natal recording of decree, the
constitutive role of intention, the causal effect of truthfulness on
barakah, and the non-depletion of wealth by charity.

\textbf{Classical kalam}: the Mu`tazili, Ash`ari, and Maturidi
positions on qadar provide the classical framework for the modern
reformulation in Chapter 16. Each position addresses a real tension;
the modern reformulation does not refute them but supplements them
with conceptual resources unavailable to the classical thinkers.

\textbf{Akbarian tradition}: from Ibn Arabi through Mulla Sadra to
Shah Wali Allah, the Akbarian system provides the philosophical
articulation of the metaphysical claims. The doctrines of
\arb{wahdat al-wujud}, divine names, tajalli, and insan al-kamil
constitute the conceptual register in which the claims receive
their most rigorous classical treatment.

\textbf{Modern scholarship}: Nasr, Chittick, Izutsu, al-Attas,
Murata, and others have made the tradition accessible to the
contemporary academic reader. Their work is the foundation on which
the present project builds.

\textbf{What comes next}: Chapter 4 explains why modern conceptual
tools are needed to formalize the claims surveyed here, and previews
how Part II will develop each of the relevant tools.
\end{summarybox}

\cleardoublepage
